У одной маленькой девочки начал гнить молочный зуб. Решили эту девочку отвести к зубному врачу, чтобы он выдернул ей ее молочный зуб.

Вот однажды стояла эта маленькая девочка в редакции, стояла она около шкапа и была вся скрюченная.

Тогда одна редакторша спросила эту девочку, почему она стоит вся скрюченная, и девочка ответила, что она стоит так потому, что боится рвать свой молочный зуб, так как должно быть, будет очень больно. А редакторша спрашивает:

<<Ты очень боишься, если тебя уколют булавкой в руку?>>

Девочка говорит:

<<Нет.>>

Редакторша уколола девочку булавкой в руку и говорит, что рвать молочный зуб не больнее этого укола. Девочка поверила и вырвала свой нездоровый молочный зуб.

Можно только отметить находчивость этой редакторши.

\medskip
\noindent\textit{6 января 1937 г.}